\section{Actual exceptional-power transport and global canonical descent}

\begin{lemma}[The whole exceptional power on the original coordinate chart]
\label{lem:kn-blowup-canonical-exceptional-power-chart}
\lean{EnvelopeProperty.Vanishing.NadelBlowupCanonicalGluingIdeals.exceptionalPowerChart_eq,
EnvelopeProperty.Vanishing.NadelBlowupCanonicalGluingIdeals.exceptionalJacobianChart_eq,
EnvelopeProperty.Vanishing.NadelBlowupCanonicalGluingIdeals.exceptionalTensorChartIso,
EnvelopeProperty.Vanishing.NadelBlowupCanonicalGluingIdeals.exceptionalTensorChartIso_inclusion}
\leanok
\uses{thm:kn-blowup-canonical-affine-ideal-iso}
On the literal native coordinate blowup, the actual exceptional ideal to
any nonnegative power restricts under the actual original model chart to
the ideal of the same coordinate power. In particular its $(r-1)$-st
power is the actual Jacobian image ideal. The genuine tensor pullback
comparison retains the entire original tensor-ideal inclusion.
\end{lemma}
\begin{proof}
Use actual geometric ideal comap multiplicativity and the native
exceptional-ideal calculation on that same chart. On the actual affine
spectrum, ideal powers of a principal ideal are the ideal of the same
power, and the genuine spectrum/global-section ring isomorphism preserves
powers. The whole native tensor-ideal open pullback formula followed by
this same ideal equality gives the full inclusion equation.
\end{proof}

\begin{lemma}[Whole native isomorphism factors from actual open charts]
\label{lem:kn-blowup-canonical-open-descent}
\lean{EnvelopeProperty.Vanishing.NadelBlowupCanonicalGluingDescent.tensorIdealMap_mono,
EnvelopeProperty.Vanishing.NadelBlowupCanonicalGluingDescent.pullback_open_mono,
EnvelopeProperty.Vanishing.NadelBlowupCanonicalGluingDescent.isIso_of_open_cover,
EnvelopeProperty.Vanishing.NadelBlowupCanonicalGluingDescent.exists_iso_of_local_factors}
\leanok
A whole original module map factors isomorphically through an original
monomorphism whenever its actual pullbacks to an actual open point-cover
have such isomorphism factors. No matching or exactness certificate is
needed.
\end{lemma}
\begin{proof}
Tensoring an ideal inclusion by an actual line is monic by genuine
invertible-tensor equivalence, and genuine open pullback retains monos.
The original map composed with the original cokernel of the inclusion
vanishes on the actual point-cover, hence vanishes on the whole scheme.
The original short exact cokernel sequence therefore produces the whole
factor. Each actual pullback of that factor equals its given local factor
by monicity. Whole monos and epis are detected on the same actual
point-cover, so this factor is an isomorphism. This derives all matching.
\end{proof}

\begin{lemma}[The native model chart retains the whole canonical factor]
\label{lem:kn-blowup-canonical-model-chart-factor}
\lean{EnvelopeProperty.Vanishing.NadelBlowupCanonicalGluingChart.exists_modelChart_factor_iso}
\leanok
\uses{thm:kn-blowup-canonical-affine-ideal-iso,lem:kn-blowup-canonical-exceptional-power-chart}
On each original native model chart, the pullback of the whole native
canonical differential has an isomorphism factor through the pullback of
the actual target canonical sheaf tensored with the actual exceptional
ideal to power $r-1$.
\end{lemma}
\begin{proof}
The genuine original model-chart triangle identifies its composite with
the blowup projection with the original chart substitution. The previous
affine producer therefore supplies its actual entire differential factor.
The actual open-immersion exterior differential identifies the pulled
target canonical sheaf with the chart canonical sheaf. The full exterior
differential composition law and actual pullback-composition isomorphism
retain the same original whole map. Tensor-ideal naturality and the
preceding actual exceptional-power chart comparison finish the factor.
\end{proof}

\begin{theorem}[The whole coordinate-blowup canonical Cartier comparison]
\label{thm:kn-blowup-canonical-whole-coordinate-iso}
\lean{EnvelopeProperty.Vanishing.NadelBlowupCanonicalGluingWhole.modelChart_point_cover,
EnvelopeProperty.Vanishing.NadelBlowupCanonicalGluingWhole.exists_coordinate_blowup_canonical_ideal_iso,
EnvelopeProperty.Vanishing.NadelBlowupCanonicalGluingWhole.coordinateBlowupCanonicalDivisorIso}
\leanok
\uses{lem:kn-blowup-canonical-model-chart-factor,lem:kn-blowup-canonical-open-descent}
For the native blowup of the original affine $n$-space over the original
field along its actual coordinate centre of codimension $m+1\leq n$,
the actual whole pulled source canonical sheaf is isomorphic to the
actual whole target canonical sheaf tensored with the actual exceptional
ideal to power $m$. Its whole inclusion composite is the original native
canonical differential. Equivalently, the actual target canonical sheaf
is the actual pulled source canonical sheaf tensored with the actual dual
of that same exceptional-power ideal, namely $\mathcal O(mE)$.
\end{theorem}
\begin{proof}
The actual original model charts cover all original blowup points. The
native smoothness theorem gives the actual target canonical line, so its
tensor with the original exceptional-power ideal includes monically.
Apply the preceding whole descent to the internally constructed original
model-chart factors. There is no local factorization, frame, cover,
matching or canonical comparison input in this original-parameter theorem.
The genuine exceptional-power Cartier module has its actual dual evaluation
isomorphism. Actual tensor associativity and this evaluation then construct
the equivalent positive-divisor formula, with the same original tensor order.
\end{proof}

\begin{remark}[Remaining boundaries]
The codimension is positive, including the genuine Cartier-centre case
$m=0$. This theorem does not assert derived direct-image acyclicity.
Transport under adapted etale charts to general smooth centres and actual
relative structure-sheaf acyclicity remain separate producers; no global
Nadel theorem or local-vanishing result is assumed here.
\end{remark}
